% Copyright (c) 2026 Eugene Joseph M. Ragasa.
% SPDX-License-Identifier: Apache-2.0
% Author: Eugene Joseph M. Ragasa
% Formal proof completed: 2026-10-19
\documentclass[11pt]{article}
\usepackage[T1]{fontenc}\usepackage{amsmath,amssymb}\usepackage[margin=1in]{geometry}\usepackage[hidelinks]{hyperref}
\title{Proof Dossier: Ragasa2026A Maximizing-Angle Range}\author{Eugene Joseph M. Ragasa}\date{Formal proof completed: 2026-10-19}
\begin{document}\maketitle
\section{Identity and status}
ID \texttt{BG-RAGASA2026A-THM-010}; status \texttt{proved}. Lean declaration:
\par\noindent {\tiny\path{ProjectKoios.BandGauge.Papers.Ragasa2026A.maximizingOperatorCorrelationAngle_mem}}.
\section{Exact statement}
The canonical complex-argument representative belongs to
$[-\pi/2,\pi/2)$.
\section{Source, proof, and boundary}
This \texttt{formalizes} the interval convention in manuscript lines 652--659 at owner commit
\begin{center}\small\texttt{99a46cb59d0c6f27ada7d709f194996ef9619d9f}.\end{center}
It uses the standard argument range
and endpoint replacement. No uniqueness or numerical cutoff decision is asserted.
Kernel axioms: \texttt{propext}, \texttt{Classical.choice}, \texttt{Quot.sound}; no
admissions. AI-assisted and kernel-checked; no independent expert review recorded.
\end{document}
